\chapter{Conclusion and Future Work}
\label{ch:conclusion}
% 개요 12장에 따라 다시 썼다. 옛 원고의 결론(대상 소속 모델링, 16실행 개발 집합과
% 127실행 재사용 집합의 소속 교환 대조, benign interleaving에서 event-wise history의
% 조건)은 옛 관계 모델 줄기의 요약이다. 그 결과는 answers 1에 따라 부록 B
% (B.3, B.4)로 옮기므로 이 장에서는 뺐다. 숫자와 인용은 바꾸지 않았고, 이 장에
% 다시 쓰지도 않았다.

PROPHET forecasts the next of 18 attacker actions from the Windows
execution and audit records received before a cutoff. The thesis
followed one line of argument through four steps. It first measured
that seeds of REL\_GLOBAL trained on the same 603 FIT runs disagree
widely once the benign worker~C appears, and that their scores on the
training conditions do not say which seed will hold up. It then traced
much of the in-house disagreement to a single file decision, whether to
COPY or to keep searching, which the seeds resolve from the observed
result for the current target in different ways. From that diagnosis
it moved to a remedy that leaves the deployed model unchanged: a
student of the same architecture follows the averaged probabilities of
five selected teachers on unlabeled environment inputs. Finally, the
remedy was tested against criteria sealed before training. In the v3
round, with 30 new seeds and 157 new runs, mean NLL on
\texttt{A+C FLAT} fell from 0.0731 to 0.0276 and all three primary
criteria passed. On ADS-24 the seed standard deviation fell from 0.163
to 0.025, and the three encoder adaptations passed their criteria on
both datasets.
% 근거: RECORD:23, 44 (측정), 104--110 (진단, 6장), 130, 137, 139 (처방, 4장),
% 166, 170 (v3), 219 (ADS-24 SD), 294 (세 인코더). 숫자는 개요대로 세 개만 다시 썼다.
% ADS-24 결과와 세 인코더의 ADS-24 쪽은 이미 열람한 평가 실행의 개발 비교다
% (8, 9장에서 밝힘). 여기서는 판정 이름만 쓰고 지위를 다시 설명하지 않았다.

These results have firm edges. CD narrows the spread across seeds; it
does not undo the cost of moving to an unseen environment. On ADS-24
most distilled models still scored worse than the majority-class
baseline at an NLL of 1.830, and only the Transformer students cleared
it. Nor does every seed gain: in each architecture--dataset combination
between 20 and 30 of the 30 seeds improved, and seeds that had been
lucky were pulled down to the level of the consensus. The diagnosis
holds for one model in one lab, and the agreement ranking used to pick
teachers did not carry over to ADS-24. What a defender gets from CD is
a forecaster whose quality depends less on the draw of the seed, at the
price of training 30 baselines before the student. It is not a
forecaster that has learned more about the new environment.
% 근거: RECORD:419--423 (환경 이동 미해결, ADS-24 기준선 아래, Transformer만 초과),
% 297 (20~30/30, 운 좋은 seed 하락), 405 (합의도 seed 선택 ADS-24 미재현),
% 427 (30개 학습 비용). 해석은 11장 sec:disc-does에서 이미 했으므로 되풀이하지 않았다.

\section{Future Work}
\label{sec:future}

One bundle of closing checks was planned before this thesis was written
and is not reported here; whatever it shows will be added to the
limitations rather than covered by new experiments. Three of the checks
require training. The first varies the teacher set, using five teachers
chosen at random, two teachers, and the full set of 30, to learn whether
training 30 baselines and ranking them by agreement is worth its cost.
The second trains students on 25\% and 50\% of the unlabeled
environment inputs to estimate how much new-environment data the KL
term needs. The third repeats the comparison with other single-model
methods on a Transformer encoder, so that the comparison does not rest
on one architecture. The remaining checks reuse frozen models. They run
inference on the 119 ADS-24 runs that no earlier stage has used, which
would give the first ADS-24 result on runs no one has looked at, and on
a fresh in-house collection that completes the method comparison
stopped at 110 runs. They also recompute the intervals with seeds
resampled along with runs, so that a verdict no longer speaks only for
these 30 seeds, and summarize temperature scaling and probability
ensembles as the strongest simple competitors.
% 근거: RECORD:488--498 (마무리 확인 묶음 표; 불리한 결과도 한계에 적고 새 실험으로
% 덮지 않는다). 결과는 쓰지 않았다(answers 추가 지시).
% TODO(사람): 개요 질문 A1에서 온도 보정·확률 앙상블 요약을 본문(11장 또는 9장)에
% 넣기로 하면 이 문단의 마지막 항목("and summarize temperature scaling ...")을 뺀다.
% 지금은 A1의 답이 없어 남겨 두었다.

A separate line of work concerns the mechanism. The account of why the
seeds diverge was built for REL\_GLOBAL in the in-house lab, and CD
passed on ADS-24 and on the general encoders without a matching
diagnosis there. Whether those models also split on one decision, and
whether their disagreement sits where the in-house lookup sits, is
open. Extending the diagnosis to ADS-24 and to the general encoders
would show whether the step from diagnosis to remedy holds outside the
setting in which it was found.
% 근거: RECORD:430--437, 6장 첫 문단 ("We did not study the mechanism separately on
% ADS-24 or for the encoders"), 11장 sec:disc-mechanism. 개요 12장 마지막 항목.

% ---------------------------------------------------------------
% 사람에게 물을 것 (12장)
% 1. 개요 질문 A1: 온도 보정·확률 앙상블 요약을 본문에 넣는지. 넣으면 Future Work
%    첫 문단의 마지막 항목을 지운다.
% 2. 옛 결론의 문장(대상 소속, 소속 교환 16/127실행, benign interleaving 조건)을
%    부록 B 첫머리 요약으로 옮길지, 그냥 버릴지.
% 3. 장 제목을 "Conclusion and Future Work"로 바꿨다. 양식상 "Conclusion"만 써야
%    하면 알려 주세요. label은 서론이 참조하는 ch:conclusion을 그대로 쓴다.
% ---------------------------------------------------------------
